\section{Propagation of alternative equations to locked diagnostics}
\label{app:alternative-propagation}

The existence of an axiomatic countermodel does not imply that it reproduces a
measured value.  We therefore propagated representative alternatives through
the same equations and applied the locked PR-III references only after
generation.  ``Correct'' in this appendix means inside the diagnostic interval
accepted by PR-III; it does not mean exact equality.

\begin{table}[H]
\centering
\small
\caption{Compact-branch propagation to the locked rubidium
$\alpha^{-1}=137.035999206(11)$ diagnostic.}
\label{tab:compact-diagnostic-propagation}
\begin{tabular}{@{}lrr@{}}
\toprule
Branch & Generated $\alpha^{-1}$ & Distance \\
\midrule
Published $(-1,\Pi_{13})$ & 137.035999207530239 & $0.1391\sigma$ \\
Alternative $(0,\Pi_{13})$ & 137.035999207470098 & $0.1336\sigma$ \\
Alternative $(+1,\Pi_{13})$ & 137.035999207409986 & $0.1282\sigma$ \\
Alternative $(-1,\Pi_{15})$ & 254.208754866216196 & fails \\
Alternative $(-1,\Pi_{17})$ & 384.885897406270487 & fails \\
\bottomrule
\end{tabular}
\end{table}

The three $\Pi_{13}$ branches are exactly distinct, but all lie inside the same
one-standard-deviation interval.  Indeed, the two nonpublished signs are
marginally closer to the diagnostic central value.  Present fine-structure
precision therefore does not select the published terminal sign.

For the PR-III electromagnetic correction, none of the 21 nonpublished
bounded-degree monomials passes the rubidium one-sigma gate when used as a
single unit-coefficient replacement term.  The published $D_1+D_2$ sum does
pass.  This is numerical discrimination only within the enumerated family; it
does not derive the exponent rule from P1 through P7. The electroweak coefficient propagation uses
\begin{equation}
 \Pi_W=\frac{a}{N_{\rm gen}}\epsilon_V,\qquad
 \Pi_Z=\frac{b}{N_{\rm gen}}\epsilon_V+
       \frac{k}{N_{\rm gen}}\epsilon_V^2.
\end{equation}
Besides the published $(a,b,k)=(2,-1,2)$, two alternatives pass both locked
one-sigma mass gates:
\begin{table}[H]
\centering
\small
\caption{Electroweak coefficient triplets passing both locked mass gates.}
\label{tab:ew-alternative-diagnostics}
\begin{tabular}{@{}lrrrr@{}}
\toprule
$(a,b,k)$ & $M_W$ (GeV) & $M_Z$ (GeV) & $M_W$ distance & $M_Z$ distance \\
\midrule
$(2,-1,1)$ & 80.37394285 & 91.18590923 & $0.3566\sigma$ & $0.8051\sigma$ \\
$(2,-1,2)$ published & 80.37394285 & 91.18767643 & $0.3566\sigma$ & $0.0364\sigma$ \\
$(2,-1,3)$ & 80.37394285 & 91.18944360 & $0.3566\sigma$ & $0.8779\sigma$ \\
\bottomrule
\end{tabular}
\end{table}

The published coefficient is closest to the $M_Z$ central value, but the
same acceptance rule does not make it unique.  Selecting it by residual would
place a diagnostic target inside the construction map and violate the no-fit
firewall.  Numerical agreement cannot repair the missing frozen-postulate
derivation.
